\documentclass[12pt]{article}
\usepackage[utf8]{inputenc}
\usepackage[T1]{fontenc}
\usepackage{hyperref}
\title{Marti Ke Language Resources for K-12 Curriculum Development}
\author{Compiled by Claude Sonnet 5.5 (claude-sonnet-5-5) for Andr\'as Kornai}
\date{}
\begin{document}
\maketitle

\section{Introduction}
Marti Ke (ISO 639-3 zmg; Glottocode mart1254), also written Mati Ke, is a Western Daly language of the Northern Territory of Australia. It is traditionally associated with the coast between the Moyle River and Port Keats (Wadeye). Sources report about 30 fluent speakers in 1983 and about 3 speakers in the early 2000s. Wikipedia lists the language as extinct by 2016, while another source describes a revival project. Speakers have mostly shifted to English and Murrinh-Patha.

\section{Primary Educational Resources}
No primers or school materials for Marti Ke were found in the searches.

\section{Linguistic and Reference Materials}
\begin{itemize}
\item Wikipedia, ``Marti Ke language''. \url{https://en.wikipedia.org/wiki/Marti_Ke_language}
\item Endangered Languages Project entry. \url{https://endangeredlanguages.com/lang/2330?hl=de}
\end{itemize}

\section{Religious and Community Resources}
\begin{itemize}
\item Wikipedia, ``Mati Ke'': about the Mati Ke people, and it mentions a language revival project. \url{https://en.wikipedia.org/wiki/Mati_Ke}
\end{itemize}

\section{Digital and Audio Resources}
No audio or digital collection was identified in the searches.

\section{Curriculum Development Context}
The search results say Marti Ke was one of the languages prioritised in the Priority Languages Support Project run by First Languages Australia (as of 2020). That is the most relevant lead for revival work. Any curriculum would depend on archival material and on the Mati Ke community. I did not find or examine the archival holdings.

\end{document}